% page 388 of the facsimile (image p-387 of the scan)
% block 154.9 x 242.0 mm ; left margin 8.0 mm
\pgc{-12.700mm}{0.000mm}{
\l{\cel{4}380}
\l{}
\l{\sou{momento}. (\sou{mekaniko})Produto di forco per disto. - DEFIS.}
\l{}
\l{monado. (\sou{filoz}.)I. Segun la doktrino da Pitagoro : unajo}
\l{\cel{3}perfekta, principo nevidebla di la kozi spiritala e ma-}
\l{\cel{3}teriala. - II.Segun Leibnitz : substanco simpla, nedi-}
\l{\cel{3}videbla, esencale agema. - III. (\sou{zool.}) Protozoo kun or-}
\l{\cel{3}gano lokomocala flogilatra, ek la klaso \textquotedbl{}monatidi\textquotedbl{}, ka-}
\l{\cel{3}rakterizata da un o du mikra flogili acesora, ye la bazo}
\l{\cel{3}di la flogilo precipua.- DEFIS.}
\l{}
\l{\sou{monadologio.}\cel{1}(\sou{filoz.}) Teorio pri la monadi.}
\l{}
\l{\sou{monako.}\cel{2}(\sou{religio katol.}) Monakulo enklostrigita qua vivas}
\l{\cel{3}for la socio ordinara. - DEFIRS.}
\l{}
\l{\sou{monarko}. Suvereno di stato qua guvernesas da chefo unika.}
\l{\cel{3}- DEFIRS.}
\l{}
\l{\sou{monarkio}. Stato guvernata da chefo unika. - DEFIRS.}
\l{}
\l{\sou{monastero}\cel{2}(\sou{biol.})\cel{2}Komenco-fazo di la mitoso (prefazo) dum}
\l{\cel{3}qua la centrosomo cirkondas su per filamenti protoplasma}
\l{\cel{3}akromata, astroforma. -}
\l{}
\l{\sou{monato}. Tempo-quanto adoptita kom un ek la dekedu dividuri}
\l{\cel{3}di yaro, e qua kontenas lore 30, lore 31 dii (ecepte la}
\l{\cel{3}monato februaro, qua kontenas 28 dii dum la yari ordinara,}
\l{\cel{3}e 29 dum la yari bisextila). - DE.}
\l{}
\l{\sou{mondo}.Ensemblo ek lo existanta. Ensemblo ek Suno qua lumizas}
\l{\cel{3}ni, Tero quan ni habitas, e la planeti cetera kun lia}
\l{\cel{3}sateliti e kometi. - EFIS.}
\l{}
\l{\sou{monero}. (\sou{biol.}) Grupo ek protozoi, kreita da Haeckel, qua}
\l{\cel{3}igis ek li klaso aparta, karakterizata per (precipue)}
\l{\cel{3}lia\cel{1}\sou{tot}a indijo di nukleo ; li esus do celuli nekompleta.}
\l{\cel{3}- DEFIS.}
\l{}
\l{\sou{moneto}. Peco de metalo qua uzesas lor kambii, kun estampuro}
\l{\cel{3}legala qua atestas lua pezo e titruro e determinas lua}
\l{\cel{3}valoro nominala. - EFIRS.}
\l{}
\l{\sou{monismo}. (\sou{filoz.}) I. Sistemo filozofiala qua explikas univer-}
\l{\cel{3}so per la existo di tipo unika di realajo. - II. Sistemo}
\l{\cel{3}da Haeckel qua expresas universo per (unike) la efiko di}
\l{\cel{3}la forci fizikala e kemiala. - DEFIRS.}
\l{}
\l{\sou{monisto}. Adepto di monismo.\cel{2}- DEFIRS.}
\l{}
\l{\sou{monitoro}. Milito-navo, di tun-konteno mezvalora, basa\cel{1}\sou{sur}\cel{1}la}
\l{\cel{3}aquo-nivelo, kurasizita tre forte, e provizita ye kanoni}
\l{\cel{3}gros-modela. - DEFIRS.}
\l{}
\l{\sou{mono-}\cel{1}.Prefixo ciencala = un.}
}
